\scanpage{69}
\mspage{13}

il suffit de définir
\[ \underbrace{iF}_{\substack{\Vert\\(p\unsure{\varphi|q})i}}\overset{\unsure{\cdot}\,\overline{v_0}}{\text{---}}i\deleted{\unsure{\mathrm{id}}} \]
qui soit dans $W'$ argument par argument, puisque $W'_{\unsure{0|o}}$ induit par $W'$ et que $F\mapsto\deleted{\unsure{F}}iF$
\[ \underline{\mathrm{Hom}}(C'_0,C'_0)\unsure{\longrightarrow|\text{---}}\underline{\mathrm{Hom}}(C'_0,C') \]
est \unsure{pl.} fid. ($i$ l'étant). Or il suffit de prendre $\overline{v_0}=v*i:(p\varphi)\circ i\ \text{---}\ \mathrm{id}_{C'}\circ i$.

Pour résumer, il faut montrer que le coup des couples de foncteurs quasi-inverses marche dans les cinq cas suivants \unsure{:|6}
\diagram{\begin{tikzcd}[ampersand replacement=\&,column sep=small,row sep=huge]
\underline{\mathrm{Fib}\,\mathrm{sc}}\arrow[drr,shift left]
 \&\underline{\mathrm{Fib}\,\mathrm{sc}}_0\arrow[l,hook']\arrow[dr,shift left]
 \&\underline{\mathrm{Parf}}\arrow[d,shift left]
 \&\underline{\mathrm{Cofib}\,\mathrm{sc}}_0\arrow[dl,shift right]
 \&\underline{\mathrm{Cofib}}\arrow[dll,shift right]\\
 \& \&\mathrm{Cat}_{/S}\arrow[ull,shift left,"\underline{\Phi}"']\arrow[ul,shift left,"\underline{\Phi}_0"']\arrow[u,shift left,"C"]\arrow[ur,shift right,"\Psi_0"']\arrow[urr,shift right,"\Psi"']\&\&
\end{tikzcd}}

D'ailleurs par dualité, on est ramené aux \inserted{trois} cas de $\underline{\mathrm{Fib}\,\mathrm{sc}}$, $\underline{\mathrm{Fib}\,\mathrm{sc}}_0$, $\underline{\mathrm{Parf}}$. (Peut être le cas $\underline{\mathrm{Fib}\,\mathrm{Sc}}_0$, avec $\underline{\Phi}_0$ défini en termes de $\underline{\Phi}$ et de $C$, se ramènerait-il aux deux cas qui l'entourent, de $\underline{\mathrm{Fib}\,\mathrm{sc}}$ et $\underline{\mathrm{Parf}}$.)

\centerline{$\unsure{\backslash}$}

Je vois que je dois rectifier le tir, tout d'abord parce que j'ai oublié que les foncteurs

% archive sheet number: 68
